\section{Proposed Method}

We use an AI-assisted mixed-methods design implemented in Python with Polars
DataFrames to analyze GitSkills~\cite{gitskills2027}, downloaded from Hugging
Face,\footnote{\url{https://huggingface.co/datasets/mvaccargiu/gitskills}}
and we fix a seed for stochastic steps (sampling and applicable model calls) so
the run is reproducible.
Fig.~\ref{fig:pipeline-simple} summarizes the five major stages, while the
appendix gives the extended flowchart (accuracy checks, manual gap fill, and
merges).

For each AI stage we apply a quality gate on a random sample of outputs (sample
size TBD), and tie expected accuracy to the operating
threshold (e.g., a GLiNER score threshold of $0.8$ requires at least $80\%$
sample accuracy).
If the sample is below target, we adjust the threshold and re-run or re-filter
that stage, and if it meets the target we manually resolve gaps, merge
automatic and manual results (dropping non-existing entities), and only then
advance.

\begin{figure}[!t]
  \centering
  \begin{tikzpicture}[
  stage/.style={
    draw,
    rounded corners=2pt,
    align=center,
    text width=0.82\columnwidth,
    inner sep=5pt,
    font=\small\sffamily
  },
  arr/.style={-{Stealth[length=2mm]}, thick}
]
  \node[stage] (s1) {Dataset population};
  \node[stage, below=5mm of s1] (s2) {Tool and organization extraction};
  \node[stage, below=5mm of s2] (s3) {Repository resolution};
  \node[stage, below=5mm of s3] (s4) {License discovery};
  \node[stage, below=5mm of s4] (s5) {License freedom grouping};
  \draw[arr] (s1) -- (s2);
  \draw[arr] (s2) -- (s3);
  \draw[arr] (s3) -- (s4);
  \draw[arr] (s4) -- (s5);
\end{tikzpicture}

  \caption{High-level analysis pipeline.
    See the appendix for the extended diagram (accuracy checks, manual gap fill,
    and merges).}
  \label{fig:pipeline-simple}
\end{figure}

\subsection{Dataset population}

We download GitSkills and load the Parquet tables with Polars, while population
filters and analysis slices remain TBD (candidates include
\texttt{dedup\_primary}, \texttt{frontmatter\_valid}, exact \texttt{SKILL.md}
filename, and \texttt{location\_class}).
For extraction we use only the \texttt{SKILL.md} content and never sibling
files, keeping a skill when that content is textual and skipping non-text
cases such as CSS.

\subsection{Tool and organization extraction}

We extract taught \emph{tools} and \emph{organizations} with GLiNER
(initial score threshold $0.8$), where a tool is a product the skill teaches
the agent to use, not the host skill directory, publishing repository, or an
in-text attribution link.
Organizations support later repository search, while RQ1 frequencies count
taught tools.
After the quality gate, empty or unresolved rows are filled manually and merged
before repository resolution.

\subsection{Repository resolution}

We resolve tools sequentially, reusing a project when a tool is already linked
and otherwise querying versioning platforms (e.g., GitHub or GitLab) and using
a reasoning model (TBD) to select the owning repository from API candidates,
with open-web search as a fallback.
The same accuracy gate, manual gap fill, and merge close this stage before
license discovery.

\subsection{License discovery}

For each resolved repository we scan the repository root for conventional
license filenames (e.g., \texttt{LICENSE}, \texttt{COPYING}), extract the
license identity, and link it to the tool, with gaps again going through sample
validation (when AI-assisted), manual review, and merge.

\subsection{License freedom grouping}

Confirmed licenses are grouped by \emph{freedom}, meaning what the license
allows or denies doing with the tool (not only SPDX\footnote{Software Package Data Exchange (SPDX) provides short standard license identifiers such as MIT and Apache-2.0. See \url{https://spdx.org/licenses/}.} string frequency), and these
groups support RQ2 (license distribution) and RQ3 (project access patterns for
top tools versus freedom class).

Units of analysis follow the pipeline, so RQ1 uses distinct skills and tool
mentions, RQ2 uses confirmed tools with licenses, and RQ3 uses skill-tool pairs
for frequent tools after merge.
Limitations include TBD sample sizes and population slices, a TBD repository
ranking model, wrong-repository selection, root-only license scans, and residual
API non-determinism where seeds are not honored.
